\documentclass[10pt,a4paper]{article}
\usepackage[margin=2.6cm]{geometry}
\usepackage[T1]{fontenc}
\usepackage{amsmath,amssymb}
\IfFileExists{newtxtext.sty}{\usepackage{newtxtext,newtxmath}}{\usepackage{txfonts}}
\usepackage[scaled=.95]{helvet}
\usepackage{microtype,booktabs,tabularx,listings,xcolor,titlesec,enumitem,needspace,placeins}
% Palette shared with Tolerant frames and An Effect Language for TScert.
\definecolor{accent}{HTML}{1F5F8B}
\definecolor{chgc}{HTML}{0B7A68}
\definecolor{regc}{HTML}{B45309}
\definecolor{asrc}{HTML}{6D3FC4}
\definecolor{rulec}{HTML}{8A1C4A}
\definecolor{codec}{HTML}{24292F}
\definecolor{ckc}{HTML}{0550AE}
\definecolor{csc}{HTML}{116329}
\definecolor{ccmc}{HTML}{6E7781}
\definecolor{leanc}{HTML}{2F6F5E}
\definecolor{conjc}{HTML}{9A6700}
\definecolor{exbg}{HTML}{F6F4EE}
\usepackage[colorlinks,linkcolor=accent,citecolor=accent,urlcolor=accent]{hyperref}

\titleformat*{\section}{\Large\bfseries\sffamily\color{accent}}
\titleformat*{\subsection}{\large\bfseries\sffamily\color{accent}}
\titleformat*{\paragraph}{\bfseries\sffamily\color{accent}}
\setlist[description]{font=\normalfont\sffamily\bfseries\color{accent}}

\lstset{basicstyle=\footnotesize\ttfamily,columns=fullflexible,keepspaces=true,
  breaklines=true,showstringspaces=false,upquote=true,xleftmargin=2em,
  numbers=left,numberstyle=\tiny\color{ccmc},numbersep=7pt,
  frame=l,rulecolor=\color{accent},framesep=6pt,backgroundcolor=\color{exbg},
  keywordstyle=\color{ckc}\bfseries,
  commentstyle=\color{ccmc},stringstyle=\color{csc}}
\lstdefinelanguage{TypeScript}{
  morekeywords={type,export,function,const,let,for,while,if,else,throw,new,return,
    readonly,number,string,void,boolean,extends,true,false,null,undefined,
    Modifies,AtMost,Effects,Join,SafeInt,Dense,
    Int32Array,Uint16Array,Uint8Array,Map},
  sensitive=true,morecomment=[l]{//},morecomment=[s]{/*}{*/},
  morestring=[b]",morestring=[b]',
}

% Core logic macros (kept in plain form here for prep.py HTML/MathJax compatibility;
% coloured for PDF typesetting via \AtBeginDocument below).
\newcommand{\sep}{\mathrel{\#}}
\newcommand{\kw}[1]{\mathsf{#1}}
\newcommand{\keep}{\kw{keep}}
\newcommand{\init}{\kw{init}}
\newcommand{\grow}{\kw{grow}}
\newcommand{\any}{\kw{any}}
\newcommand{\old}{\kw{old}}
\newcommand{\val}{\kw{values}}
\newcommand{\st}{\kw{spines}}
\newcommand{\fld}{\kw{fields}}
\newcommand{\len}[1]{|#1|}
\newcommand{\code}[1]{\textcolor{codec}{\texttt{#1}}}
\newcommand{\lean}[1]{\texttt{\small #1}}
\newcommand{\ruleName}[1]{\textcolor{rulec}{\textsc{#1}}}
\newcommand{\chg}[1]{\textcolor{chgc}{\mathsf{#1}}}
\newcommand{\reg}[1]{\textcolor{regc}{\mathsf{#1}}}
\newcommand{\asr}[1]{\textcolor{asrc}{\mathsf{#1}}}
\newcommand{\reach}[2]{#1.{\ast\ast}_{#2}}
\newcommand{\modstate}[1]{#1.{@\kw{state}}}
\newcommand{\status}[1]{\textcolor{leanc}{\textsf{\small[#1]}}}
\newcommand{\statusc}[1]{\textcolor{conjc}{\textsf{\small[#1]}}}
\newcommand{\statuso}[1]{\textcolor{accent}{\textsf{\small[#1]}}}
\newcommand{\nI}[1]{\emph{Tolerant frames}~\cite{tf}, \emph{#1}}
\newcommand{\entry}[1]{\paragraph{#1.}}

\AtBeginDocument{%
  \renewcommand{\sep}{\mathrel{\textcolor{asrc}{\#}}}%
  \renewcommand{\keep}{\chg{keep}}%
  \renewcommand{\init}{\chg{init}}%
  \renewcommand{\grow}{\chg{grow}}%
  \renewcommand{\any}{\chg{any}}%
  \renewcommand{\old}{\asr{old}}%
  \renewcommand{\val}{\reg{values}}%
  \renewcommand{\st}{\reg{spines}}%
  \renewcommand{\fld}{\reg{fields}}%
  \renewcommand{\lean}[1]{\textcolor{leanc}{\texttt{\small #1}}}%
}

\frenchspacing
\emergencystretch=3em

\title{\sffamily\bfseries\color{accent}Tolerant Frames in Practice: Pretext's Patterns\\[2pt]
  \large\normalfont\color{black}An Empirical Catalogue of Bound-Proof Obligations in a Text-Layout Library}
\author{TScert Working Paper}
\date{27 November 2026}

\begin{document}
\maketitle

\begin{abstract}\noindent
This paper catalogues the index-bound obligations extracted from the
TypeScript text-layout library Pretext~\cite{pretext} to guide and
evaluate TScert's framing, effect, and arithmetic rules. It accompanies
\emph{Tolerant frames}~\cite{tf} and \emph{An Effect Language for
TScert}~\cite{el}.

In Pretext's baseline bound-audit corpus (\code{9268edb1}), 422 index
reads carry runtime fallbacks or \code{RangeError} guards because the
baseline pass finds no proof of $0\le k<\len{x}$. At 215 of these sites,
the local proof requires a rule beyond the baseline pass (\emph{blocked}
sites); 204 sites need only caller preconditions or builder
postconditions; and 3 sites fall outside any pattern.

The 215 blocked sites reduce to 21 minimal patterns P1--P21, together
with P22, which arises at P1's higher-order walker sites once visitors
state effects. Seven patterns---P1, P2, P3, P8, P12, P17 (118 sites; 48
of P1's walker sites additionally require P22 under higher-order
callbacks), and standalone P21 (1 site behind P5)---verify on TScert's
proved single-pass verifier (\lean{SinglePass}), alongside the
fresh-path separation half of P4 and fresh-view postfix fills of P13.
Nine mechanized rules in \emph{An Effect Language for TScert}~\cite{el}
address P22, P6's dynamic field writes, and P7's visitor closures; in
its prototype mode on upstream Pretext (\code{0dafecb3}), all 105/105
units across \code{line-break.ts}, \code{layout.ts},
\code{rich-inline.ts}, and \code{prepare.ts} verify with 0 unproved
assumptions and all 56 walker/stepper guards restored, while uncovering
2 upstream bounds bugs and 2 public-API cursor boundary checks. Each
catalogue entry states its minimal program, proof obligation,
verification principle, and verifier status.
\end{abstract}

\section*{Scope and Epistemic Status}

This paper is an empirical specification and design benchmark for
TScert's bound-verification rules. Claims fall into four epistemic
tiers:
\begin{itemize}[leftmargin=1.5em,itemsep=2pt]
\item \textbf{Observed on the corpus:} Site counts come from the
  baseline Pretext bound-audit corpus (\code{9268edb1}), and end-to-end
  walker and builder verification counts come from upstream Pretext
  (\code{0dafecb3}); empirical measurements carry \statuso{observed}.
\item \textbf{Proved in Lean:} Patterns discharged by TScert's proved
  single-pass verifier carry \status{Verified in SinglePass} and cite
  their Lean theorem names. Model theorems of \emph{Tolerant
  frames}~\cite{tf} and the nine-rule idealised heap model of \emph{An
  Effect Language for TScert}~\cite{el} (\lean{model\_rules}) are cited
  by rule and theorem name.
\item \textbf{Prototype-evaluated and conjectured:} Patterns discharged
  in the nine-rule model and exercised through \code{paper-rules-prototype}
  carry \statuso{Model + prototype}; full-calculus soundness of those
  rules against TScert's operational semantics is conjectured
  \statusc{conj.}. Open capability gaps carry \statusc{Open extension}.
\item \textbf{Assumed at external boundaries:} Properties of unanalysed
  host libraries, offline table generators, and external API callers
  cannot be proved from the analysed source and form explicit trust
  obligations (Section~\ref{sec:trust}).
\end{itemize}

\noindent
Three specification mechanisms appear throughout the catalogue:
\emph{footprint and contract annotations} (\code{Modifies<P, T>} and
\code{Effects<P, T>} assign relational change bounds $\keep$, $\init$,
$\grow$, $\any$ to heap regions such as $\val$, $\st$, $\fld$, or
parameter paths \code{"out[*]"} and \code{"c.value"}, alongside
\code{@requires}, two-state \code{@ensures} with \code{old(...)},
\code{@invariant}, and \code{@stable} field declarations);
\emph{checked local-object confinement} (an object allocated locally
that never escapes into a callee's operands retains facts supported
solely by its own locations across arbitrary-effect calls, proved by
\lean{TolerantFramesConfined.run\_sound\_pass}); and \emph{reviewed host
contracts} (audited length and effect postconditions for built-in
JavaScript methods).

\section{Corpus, Methodology, and Flagship Verified Patterns}\label{sec:corpus}

\subsection{Bound Rewrite Sites, Obligations, and Corpus Pins}

Under TypeScript's \code{noUncheckedIndexedAccess}, an index read
\code{x[k]} has type \code{T | undefined}. TScert treats a non-null
assertion \code{x[k]!} not as an unchecked cast, but as a checked proof
obligation requiring $k:\mathit{SafeInt}$, $0\le k<\len{x}$, and element
presence (density) at index $k$. Here $\mathit{SafeInt}$ denotes an
IEEE-754 binary64 integer with $|k|\le M=2^{53}-1$ (written
\code{safeInteger(k)} in annotations), the range in which integer
addition and subtraction are exact whenever the result stays in
$[-M,M]$.

We pin two Pretext commits~\cite{pretext} for distinct roles: commit
\code{9268edb1} is the baseline rewritten audit corpus from which the
422-site taxonomy and patterns P1--P21 are extracted, whereas upstream
commit \code{0dafecb3} is the target of the end-to-end walker and
builder verification of \emph{An Effect Language for TScert}~\cite{el}
(Section~\ref{sec:p22}). In the baseline corpus, a \emph{bound rewrite
site} is an array, typed-array, or string index read \code{x[k]}
rewritten as \code{x[k] ?? missing()} or guarded by a \code{RangeError}
test because the baseline verifier finds no proof of $0\le k<\len{x}$.
A site is \emph{blocked} if proving its bound locally requires a
deduction rule or specification form beyond the baseline verifier;
otherwise it is \emph{caller/builder-establishable}. A \emph{pattern} is
the smallest standalone TypeScript program isolating the obligation at
its group of blocked sites.

\begin{center}\small
\begin{tabularx}{\textwidth}{@{} >{\raggedright\arraybackslash}p{0.35\textwidth} r >{\raggedright\arraybackslash}X@{}}
\toprule
Category & Sites & Role in this catalogue\\
\midrule
Blocked --- verified in \lean{SinglePass} (P1, P2, P3, P8, P12, P17) & 118 & P1, P2, P8 in Section~\ref{sec:verified} (48 of P1's walker sites also meet P22); P3, P12, P17 illustrated in Sections~\ref{sec:framing} and~\ref{sec:arith}.\\
Blocked --- partial, model/prototype, or open (P4--P7, P9--P11, P13--P16, P18--P20) & 97 & Sections~\ref{sec:framing}--\ref{sec:external} (P21's 1 standalone-verified site sits behind P5; P4's \code{paths} and P13's fresh-view fill verify in \lean{SinglePass}).\\
Caller/builder-establishable & 204 & Require stronger caller preconditions or builder postconditions, not new rules.\\
Unclassified & 3 & Characterized individually in Section~\ref{sec:summary} (34 non-index markers excluded).\\
\midrule
\textbf{Total bound rewrite sites} & \textbf{422} & Pretext baseline bound-audit corpus (\code{9268edb1}).\\
\bottomrule
\end{tabularx}
\end{center}

\subsection{Flagship Patterns Verified by TScert's Single Pass}\label{sec:verified}

Four foundational patterns illustrate how Tolerant Frames verifies
idiomatic mutable TypeScript directly in TScert's proved single pass
(\lean{SinglePass}), without ownership types.

\paragraph{P1: Visitor with closed growth effect (50 sites; 48 walker sites also meet P22).}
\status{Verified in SinglePass}

\begin{lstlisting}[language=TypeScript]
type BoundedVisitor = (w: number) => Modifies<{ values: "any"; spines: "grow" }, void>;

function walkBounded(xs: number[], visit: BoundedVisitor | undefined): number {
  let s = 0;
  for (let i = 0; i < xs.length; i++) {
    s += xs[i]!;
    if (visit !== undefined) visit(s);
  }
  return s;
}
\end{lstlisting}

\noindent
Stating the callback's closed growth effect
$\{\val:\any,\ \st:\grow\}$ lets axiom \ruleName{(Grow)} of
\emph{Tolerant frames}~\cite{tf} preserve the lower bound
$i<\len{\mathit{xs}}$ across \code{visit(s)} even when the visitor may
append to \code{xs}. Standalone P1 verifies directly on
\lean{SinglePass}. In Pretext, however, 48 of P1's 50 sites lie inside
higher-order line walkers whose callers mutate their own arrays and
whose prepared handles maintain cross-array length equalities such as
$\len{\mathit{segmentFlags}}=\len{\mathit{widths}}$; a closed
$\st:\grow$ effect drops such equalities and blocks caller writes, so
those 48 walker sites additionally require P22 (Section~\ref{sec:p22})
and verify in \code{paper-rules-prototype}.

\needspace{16\baselineskip}
\paragraph{P2: Growth by a known count across a callee (31 sites).}
\status{Verified in SinglePass}

\begin{lstlisting}[language=TypeScript]
type Lists = { /** @stable */ readonly widths: number[] };

/** @requires safeInteger(n) && 0 <= n && n <= 1000 && 0 <= lists.widths.length
 *  @ensures lists.widths.length === old(lists.widths.length) + n */
function measure(lists: Lists, n: number):
  Modifies<{ values: "any"; spines: "keep"; "lists.widths[*]": "any" }, void> {
  /** @invariant safeInteger(k) && safeInteger(n) && 0 <= k && k <= n
   *             && lists.widths.length === old(lists.widths.length) + k */
  for (let k = 0; k < n; k++) lists.widths.push(k);
}
\end{lstlisting}

\noindent
When a caller invokes \code{measure(lists, 10)} inside a loop over a
typed array \code{flags} of length \code{cap}, the two-state
postcondition $\len{\mathit{lists}.\mathit{widths}} =
\old(\len{\mathit{lists}.\mathit{widths}}) + n$ instantiates at the
call's entry state to bound \code{lists.widths.length - 1} below
\code{flags.length}, while the regional effect
\code{"lists.widths[*]": "any"} with $\st:\keep$ preserves
$\len{\mathit{flags}}$.

\paragraph{P8: Length and presence after a write at the length (10 sites).}
\status{Verified in SinglePass}

\begin{lstlisting}[language=TypeScript]
/** @requires safeInteger(k) && 0 <= k && k <= xs.length && k <= 1000000 */
export function readWritten(xs: number[], k: number): number {
  xs[k] = 1;
  return xs[k]!;
}
\end{lstlisting}

\noindent
Writing at \code{xs[k]} with $0\le k\le\len{\mathit{xs}}$ satisfies
$\len{\mathit{xs}} = \max(\old(\len{\mathit{xs}}), k+1)$ via rules
\ruleName{Write-In} and \ruleName{Write-End} of \emph{Tolerant
frames}~\cite{tf} (proved by \lean{atLengthAtoms\_sound}), covering both
direct reads at the written index and builders that append via
\code{widths[at] = mi}. Note that in \emph{Tolerant frames}~\cite{tf},
\ruleName{Write-End} specifies $\len{x}=\old(\len{x})+1 \wedge
x[\old(\len{x})]\doteq v$ (leaving standalone $\old(\mathsf{dense}(x,i))
\Rightarrow \mathsf{dense}(x,i+1)$ under future work); typing \code{xs[k]!} on
\code{number[]} in \lean{SinglePass} relies on its mechanized
dense-write rule (\lean{denseAppendAtoms\_sound}).

\paragraph{P21: Shift-left and bitwise-OR table indexing (1 site, behind P5).}
\status{Verified in SinglePass}
The two-stage table read \code{values[(hi <{}< 8) | (c \& 0xff)]!}
verifies in \lean{SinglePass} via rules \ruleName{Shr}, \ruleName{Shl},
\ruleName{Mask}, and \ruleName{Or} (\nI{Arithmetic steps};
\lean{ArithmeticRules.rules\_sound}). In Pretext's \code{getClass}, this
site sits behind P5 (Section~\ref{sec:p5}), which bounds the decoded
table entry \code{hi}.

\medskip\noindent
Each entry in Sections~\ref{sec:framing}--\ref{sec:external} presents
four uniform parts: a minimal TypeScript program, its \emph{Obligation}
(the bound or invariant to discharge), its \emph{Verification principle}
(the logical rule that discharges it), and its \emph{Verifier status}
across \lean{SinglePass}, \emph{Tolerant frames}~\cite{tf}, and
\emph{An Effect Language for TScert}~\cite{el}. Patterns P3, P12, and
P17 remain illustrated in full as positive demonstrations of Tolerant
Frames in \lean{SinglePass}, cleanly distinguished from open extensions.

\needspace{28\baselineskip}
% ======================================================================
\section{Framing, Aliasing, and Callback Patterns}\label{sec:framing}
% ======================================================================

Five patterns (51 sites plus P22) concern preserving heap facts across
writes and calls: separating heap location kinds (P3, P6), separating
fresh allocations from field paths and checking index comparisons (P4),
framing multi-field rigid guards (P18), and higher-order callbacks
(P22).

\needspace{20\baselineskip}
\subsection{P3: Quantified Fact Across a Write Outside Its Array (19 sites)}\label{sec:p3}

\begin{lstlisting}[language=TypeScript]
type Fragment = { item: number; width: number };
/** @requires forall i in items: safeInteger(items[i])
 *            && 0 <= items[i] && items[i] < widths.length */
export function acrossWrite(items: Int32Array, widths: number[]): number {
  const fragment: Fragment = { item: 0, width: 0 };
  let s = 0;
  for (let i = 0; i < items.length; i++) {
    fragment.width = s;
    s += widths[items[i]!]!;
  }
  return s + fragment.width;
}
\end{lstlisting}

\entry{Obligation}
Prove $0\le\mathit{items}[i]<\len{\mathit{widths}}$ at
\code{widths[items[i]!]!} (line~9) from the quantified precondition
$\forall j<\len{\mathit{items}}.\ 0\le\mathit{items}[j]<\len{\mathit{widths}}$,
preserving that quantified fact across the object field write
\code{fragment.width = s} (line~8).

\entry{Verification principle}
Location-kind disjointness under Theorem ``Tolerant frame'' of
\emph{Tolerant frames}~\cite{tf}. Supporting a quantified fact over an
array or typed array by its element cells and spine places its support
in the element-cell and spine location kinds, which are disjoint by kind
from object field locations regardless of aliasing. Because
\code{fragment.width} is an object field (and \code{fragment} is locally
allocated), its write respects the quantified fact's support and leaves
the fact intact.

\entry{Verifier status}
\status{Verified in SinglePass} Both \code{acrossWrite} and the variant
repeating the quantified fact as a loop invariant verify in
\lean{SinglePass}. Quantified support (\lean{Survival.everyCells})
places array elements in \lean{Region.pathCells} (\lean{KindSet.elems}),
which \lean{RegionBounds.boundQuery} separates by location kind from
\lean{writeFrame}'s field location (\lean{KindSet.fieldsAndCells}) in
\lean{HeapExpressionRules.frameKept}, reinforced by local confinement.

\needspace{24\baselineskip}
\subsection{P4: Upper Bound of a Length Across a Call Growing Spines (18 sites)}\label{sec:p4}

\begin{lstlisting}[language=TypeScript]
export type Analysis = {
  /** @stable */ readonly starts: number[];
  /** @stable */ readonly flags: number[];
};
type Lists = { /** @stable */ readonly widths: number[] };

function measure(n: number, p: Lists | null):
  Modifies<{ values: "any"; spines: "keep"; "p.widths[*]": "any" }, number> {
  if (p !== null) p.widths.push(n);
  return n;
}

/** @requires a.flags.length === a.starts.length */
export function paths(a: Analysis): number {
  const widths: number[] = [];
  const lists: Lists = { widths };
  let s = 0;
  /** @invariant safeInteger(i) && 0 <= i && a.flags.length === a.starts.length
   *             && a.flags !== lists.widths && a.starts !== lists.widths */
  for (let i = 0; i < a.flags.length; i++) {
    measure(i, lists);
    s += a.starts[i]!;
  }
  return s;
}

/** @requires flags.length === starts.length && flags !== out && starts !== out
 *            && safeInteger(i) && 0 <= i */
export function straight(
  flags: number[], starts: number[], out: number[], i: number,
): number {
  if (i + 1 < flags.length) {
    noteIn(out, i);
    return starts[i + 1]!;
  }
  return 0;
}
\end{lstlisting}

\noindent
Here \code{noteIn(out, v)} pushes \code{v} onto \code{out} with effect
\code{Modifies<\{ values: "any"; spines: "keep"; "out[*]": "any" \}, void>}.

\entry{Obligation}
Two obligations arise when an index upper bound must survive a call
growing another spine. In \code{paths}, the loop must establish
$a.\mathit{flags}\sep\mathit{lists}.\mathit{widths}$ and
$a.\mathit{starts}\sep\mathit{lists}.\mathit{widths}$ so that
\code{measure(i, lists)}'s regional spine effect
\code{"p.widths[*]": "any"} preserves
$\len{a.\mathit{flags}}=\len{a.\mathit{starts}}$. In \code{straight},
the guard \code{i + 1 < flags.length} (line~32) must establish both
exactness of \code{i + 1} and $i+1<\len{\mathit{flags}}$ so
\code{starts[i + 1]!} (line~34) is in bounds after \code{noteIn(out, i)}.

\entry{Verification principle}
Two rules discharge these obligations:
\begin{enumerate}[leftmargin=1.5em,itemsep=2pt]
\item \emph{Fresh-array separation from field paths:} By rule
  \ruleName{Alloc} of \emph{Tolerant frames}~\cite{tf} (proved for empty
  arrays by \lean{alloc\_post}), a freshly allocated array is apart from
  every array reachable through fields of pre-existing variables,
  establishing $a.\mathit{flags}\sep\mathit{lists}.\mathit{widths}$ and
  $a.\mathit{starts}\sep\mathit{lists}.\mathit{widths}$ at allocation.
\item \emph{Exactness from a length comparison:} For $i:\mathit{SafeInt}$
  and any length $n\le M$ ($M=2^{53}-1$), a successful JavaScript
  comparison \code{i + 1 < n} implies $i<M$ (since at $i=M$, binary64
  rounding gives $\mathrm{rn}(i+1)=2^{53}>M\ge n$), hence exactness
  $|i+1|\le M$ and $i+1<n$. The $i:\mathit{SafeInt}$ premise is
  load-bearing (for a general binary64 float such as $i=-10^{20}$,
  $\mathrm{rn}(i+1)=i<n$ holds without exactness); once
  $i:\mathit{SafeInt}$ holds, monotonicity of binary64 rounding and
  exactness of $2^{53}$ cover every array, typed-array, and string
  length $n\le M$.
\end{enumerate}

\entry{Verifier status}
\status{Verified in SinglePass} for \code{paths}: \lean{SinglePass}
establishes separation between fresh literal arrays and parameter field
paths via \lean{IndexFact.apart} (\lean{PathRegionsApart}) and
\lean{restateSpines}. \statusc{Open extension} for \code{straight}:
\code{safeInteger(i)} entails only $|i|\le M$ (not $i<M$), so at $i=M$
the sum \code{i + 1} evaluates to $2^{53}$ outside $\mathit{SafeInt}$,
blocking \ruleName{Add}; neither \lean{SinglePass} nor the companion
papers yet include the safe-integer comparison rule above.

\needspace{22\baselineskip}
\subsection{P6: Quantified Fact Across a Call Mutating Value Fields (12 sites)}\label{sec:p6}

\begin{lstlisting}[language=TypeScript]
type Cache = { value: number | null };

function touch(c: Cache): Modifies<{ values: "any"; spines: "keep" }, void> {
  c.value = null;
}

/** @requires forall i in starts: safeInteger(starts[i])
 *            && 0 <= starts[i] && starts[i] < offsets.length */
export function acrossTouch(
  starts: number[], offsets: Int32Array, c: Cache,
): number {
  let s = 0;
  /** @invariant safeInteger(k) && 0 <= k
   *  @invariant forall i in starts: safeInteger(starts[i])
   *             && 0 <= starts[i] && starts[i] < offsets.length */
  for (let k = 0; k < starts.length; k++) {
    touch(c);
    s += offsets[starts[k]!]!;
  }
  return s;
}
\end{lstlisting}

\entry{Obligation}
Preserve the quantified invariant
$\forall i<\len{\mathit{starts}}.\ 0\le\mathit{starts}[i]<\len{\mathit{offsets}}$
across \code{touch(c)} (line~17) so \code{offsets[starts[k]!]!}
(line~18) is in bounds. With the coarse effect
$\{\val:\any,\ \st:\keep\}$, \code{touch(c)} permits modifying any
value-holding location, including the element cells of \code{starts}.

\entry{Verification principle}
Refine the callee's effect to object fields, separating field writes
from quantified array element cells by location kind. When the callee
mutates a field reached from a parameter path (as in \code{touch(c)}), a
parameter-rooted field region (\code{"c.value": "any"} or
\code{"c.*": "scalar"}) suffices. At Pretext's 12 sites (in
\code{han-kerning.ts} and \code{prepare.ts}), however, helpers mutate
fields of cached or dynamically selected objects not rooted at any fixed
parameter path; framing quantified array facts across those calls
requires the field-kind class region $\{\fld:\any\}$ (surface syntax
\code{fields: "any"}), separated from array element cells by rule
\ruleName{Fields-Disj} of \emph{An Effect Language for TScert}~\cite{el}.

\entry{Verifier status}
\status{Verified in SinglePass} for parameter-rooted field regions:
\lean{SinglePass} certifies parameter field regions
(\lean{ArgumentSpineRegionsWithin}) and separates them from
\lean{Region.pathCells} in \lean{PromiseCalls.effectKept}, so
\code{acrossTouch} verifies once \code{touch(c)} declares
\code{Modifies<\{ "c.value": "any" \}, void>} (or a monotone
\code{values: "init"} effect). \statuso{Model + prototype} for
Pretext's 12 dynamic-field sites: \ruleName{Fields-Disj} is mechanized
in \emph{An Effect Language for TScert}~\cite{el} (\lean{model\_rules})
and exercised in \code{paper-rules-prototype}, while \lean{SinglePass}'s
\code{Modifies} syntax does not yet include the \code{fields} class
region.

\needspace{24\baselineskip}
\subsection{P18: Guarded Fact Through Two \texorpdfstring{\code{@stable}}{@stable} Fields (2 sites)}\label{sec:p18}

\begin{lstlisting}[language=TypeScript]
type Items = { /** @stable */ readonly insideExtras: number[] | null };
/** @invariant this.items === null || this.items.insideExtras === null
 *             || this.items.insideExtras.length === this.widths.length */
export type Prepared = {
  /** @stable */ readonly widths: number[];
  /** @stable */ readonly items: Items | null;
};
type Cursor = { segment: number };

function step(c: Cursor): Modifies<{ values: "any" }, void> {
  c.segment++;
}

export function acrossStep(p: Prepared, c: Cursor): number {
  let s = 0;
  /** @invariant safeInteger(i) && 0 <= i */
  for (let i = 0; i < p.widths.length; i++) {
    step(c);
    const items = p.items;
    if (items !== null) {
      const extras = items.insideExtras;
      if (extras !== null) s += extras[i]!;
    }
  }
  return s;
}
\end{lstlisting}

\entry{Obligation}
Prove $i<\len{\mathit{extras}}$ at \code{extras[i]!} (line~22) from
\code{Prepared}'s type invariant, a guarded length equality whose guard
traverses two \code{@stable} fields (\code{p.items} and
\code{items.insideExtras}), preserved across \code{step(c)} (line~18)
under effect $\{\val:\any\}$.

\entry{Verification principle}
Axiom \ruleName{(Rigid)} of \emph{Tolerant frames}~\cite{tf}: a fact
preserved of a variable holds along any path of rigid (\code{@stable})
fields whenever the step leaves the relevant spines unchanged (or
touches only unrelated field keys). Because both \code{items} and
\code{insideExtras} are \code{@stable} and \code{step} mutates no
spines, the two-field guarded invariant survives the call.

\entry{Verifier status}
\statusc{Model only / open in SinglePass} \lean{SinglePass} verifies
single-field guarded invariants across calls and preserves multi-field
guarded facts across direct sibling field writes (\lean{writeFrame}),
but a call carrying $\val:\any$ still drops a two-field guard.

\needspace{22\baselineskip}
\subsection{P22: Local State and Shared Handles Across a Higher-Order Callback (1 call-site pattern; 48 walker sites)}\label{sec:p22}

\begin{lstlisting}[language=TypeScript]
type Visitor = (w: number) => Modifies<{ values: "any"; spines: "grow" }, void>;
type Cursor = { i: number };

export function walk(xs: number[], visit: Visitor): number {
  const c: Cursor = { i: 0 };
  let s = 0;
  /** @invariant safeInteger(c.i) && 0 <= c.i */
  while (c.i < xs.length) {
    visit(c.i);
    s += xs[c.i]!;
    c.i++;
  }
  return s;
}
\end{lstlisting}

\entry{Obligation}
Preserve both walker-local state and shared handle invariants across a
higher-order callback invocation (\code{visit(c.i)} on line~9, or
\code{onLine} in Pretext's line walkers) while allowing each caller's
callback to mutate its own state. P22 separates two distinct obligations
in Pretext's line walkers (\code{line-break.ts}):
\begin{enumerate}[leftmargin=1.5em,itemsep=2pt]
\item \emph{Form (a) --- Local cursor (unreachability):} A cursor
  allocated inside the walker and never exposed to the callback (such as
  \code{c} in \code{walk}). A closed callback effect naming
  $\val:\any$ covers $c.i$ and drops $c.i<\len{\mathit{xs}}$ even though
  the callback cannot reach \code{c}.
\item \emph{Form (b) --- Shared handle (bounded read-only reach):} A
  prepared handle passed to helpers after \code{onLine}. Unlike the
  local cursor, the handle is \emph{not} unreachable: callers hold
  \code{prepared} in scope and pass it into \code{onLine} closures (for
  instance, \code{layoutWithLines} calls
  \code{createLayoutLine(prepared, ...)} inside \code{onLine}). In
  \code{walkPreparedLinesRaw}, every iteration invokes
  \code{stepPreparedSimpleLineGeometry} on \code{prepared} and
  \code{cursor} after \code{onLine}; a closed $\st:\grow$ callback
  effect drops handle length equalities such as
  $\len{\mathit{segmentFlags}}=\len{\mathit{widths}}$ and rejects
  callers that append to their own arrays.
\end{enumerate}

\entry{Verification principle}
Bounded effect polymorphism in \emph{An Effect Language for
TScert}~\cite{el}. Define the named effect preserving a line-break
handle's declared reach (\cite{el}, Section~6):
\[
  \mathit{Geom}(p)\;\triangleq\;\{\reach{p}{\mathit{PreparedLineData}}:\keep;\ \kw{else}:\any\}.
\]
The walker quantifies over the callback's effect variable $\varepsilon$
bounded by $\varepsilon\sqsubseteq\mathit{Geom}(\mathit{prepared})$ (or
$\varepsilon\sqsubseteq\{xs.\kw{spine}:\keep;\ \kw{else}:\any\}$ for
\code{walk}). Across the callback, the two forms are discharged by
complementary rules:
\begin{itemize}[leftmargin=1.5em,itemsep=2pt]
\item \emph{Form (a)} uses \ruleName{Anch}, bounding $\varepsilon$ by
  $\{\old:\any\}$ so that locations allocated after walker entry (such
  as \code{c.i} or \code{cursor}) survive by unreachability.
\item \emph{Form (b)} uses \ruleName{Bound}, preserving everything in
  $\mathit{Geom}(\mathit{prepared})$ because the callback's effect bound
  forbids mutating the handle's declared reach even when the callback
  reads \code{prepared}.
\end{itemize}
Rule \ruleName{Meet} combines both bounds to preserve mixed relations
such as $c.i<\len{\mathit{xs}}$ across the callback, and rule
\ruleName{Inst} discharges each caller's concrete effect instantiation
(paired with \ruleName{Confined} for callbacks that populate
module-private caches).

\entry{Verifier status}
For Form (a) alone, \lean{SinglePass} preserves facts supported
exclusively by confined local objects (such as \code{safeInteger(c.i)}
and \code{0 <= c.i}) across arbitrary-effect callbacks via checked
local-object confinement (\lean{TolerantFramesConfined.run\_sound\_pass}),
whereas mixed cursor/spine relations and Form (b) require effect
polymorphism. \statuso{Model + prototype} All nine rules of \emph{An
Effect Language for TScert}~\cite{el} (\ruleName{Bound},
\ruleName{Anch}, \ruleName{Meet}, \ruleName{Inst}, \ruleName{Confined},
\ruleName{Callback-Inv}, \ruleName{Fields-Disj},
\ruleName{Plain-Cells-Disj}, and \ruleName{Elem-Len-Extend}) are
mechanized in Lean over an idealised heap model
(\lean{Verification.Paper.EffectLanguage.Model.model\_rules});
full-calculus soundness against TScert's operational semantics is
conjectured \statusc{conj.}. Exercised through
\code{paper-rules-prototype} on upstream Pretext (\code{0dafecb3}), all
105/105 units across \code{line-break.ts} (17), \code{layout.ts} (18),
\code{rich-inline.ts} (38), and \code{prepare.ts} (32) verify with 0
\code{paperAssume}s and all 56 walker and stepper bound guards restored
\statuso{observed}. That verification also uncovered 2 genuine upstream
Pretext bounds bugs (\code{getGraphemeStart} in \code{line-text.ts},
where an out-of-range grapheme index read \code{undefined} typed as
\code{number} via \code{!}, yielding full-segment slices and \code{NaN}
source offsets; and \code{getCanvasFont} in \code{measurement.ts}, where
indexing past a short generic-family list inserted \code{"undefined"}
into CSS font strings) and prompted 2 public-API cursor boundary checks
at streaming entry points (\code{stepNextLine} in \code{layout.ts} and
\code{layoutNextRichInlineLineRange} in \code{rich-inline.ts})
validating caller-supplied cursors.

% ======================================================================
\section{Closure, Module, and Collection Invariants}\label{sec:state}
% ======================================================================

Four patterns (19 sites) rely on invariants established outside the
immediate function body: facts over a local closure's captures (P7),
module-scoped constant arrays and lazy caches (P15, P16), and values
retrieved from a \code{Map} (P20).

\needspace{22\baselineskip}
\subsection{P7: Fact Over a Local Closure's Captures (10 sites)}\label{sec:p7}

\begin{lstlisting}[language=TypeScript]
type Glyph = { advance: number };

export function groups(): number {
  const glyphs: Glyph[] = [
    { advance: 1 }, { advance: 2 }, { advance: 3 }, { advance: 4 },
  ];
  /** @requires safeInteger(from) && safeInteger(to)
   *            && 0 <= from && to <= glyphs.length */
  const group = (from: number, to: number): number => {
    let s = 0;
    /** @invariant safeInteger(i) && safeInteger(to)
     *             && 0 <= i && to <= glyphs.length */
    for (let i = from; i < to; i++) s += glyphs[i]!.advance;
    return s;
  };
  return group(0, 2) + group(2, 4);
}
\end{lstlisting}

\entry{Obligation}
Establish the loop invariant $0\le i$ and
$\mathit{to}\le\len{\mathit{glyphs}}$ inside the local closure
\code{group} (lines~11--12), and discharge its precondition
\code{to <= glyphs.length} at the calls \code{group(0, 2)} and
\code{group(2, 4)} using the enclosing fact $\len{\mathit{glyphs}}=4$.

\entry{Verification principle}
Two related mechanisms govern closures over immutable captures. For
non-escaping visitor closures that read enclosing invariants without
parameter-to-capture preconditions (such as Pretext's visitor closures
in \code{rich-inline.ts}), rule \ruleName{Callback-Inv} of \emph{An
Effect Language for TScert}~\cite{el} carries facts over immutable
captures into the callback body whenever the host walker preserves their
support. For local helper closures whose parameters are bounded against
a captured array (such as \code{group} or Pretext's
\code{getHanKerningTrims}), \code{@requires} must also permit clauses
naming immutable captures, discharged against the enclosing state at
each call site.

\entry{Verifier status}
\statuso{Model + prototype (visitors); open (local helpers)}
\lean{SinglePass} verifies \code{@requires} on local arrow helpers when
clauses name parameters alone (and verifies the same loop directly in
the enclosing body), but rejects capture-referencing \code{@requires}.
\ruleName{Callback-Inv} is mechanized in \emph{An Effect Language for
TScert}~\cite{el} (\lean{model\_rules}) and exercised in
\code{paper-rules-prototype} on Pretext's visitor closures;
capture-referencing \code{@requires} on general local helpers remains an
open extension.

\needspace{18\baselineskip}
\subsection{P15: Module Constant Array Length at Entry (4 sites)}\label{sec:p15}

\begin{lstlisting}[language=TypeScript]
const KINDS: readonly string[] = ["a", "b", "c", "d"];

/** @requires safeInteger(f) && 0 <= f && f < 256 */
export function kindName(f: number): string {
  return KINDS[f & 3]!;
}
\end{lstlisting}

\entry{Obligation}
Prove $f\mathbin{\&}3<\len{\mathit{KINDS}}$ at \code{KINDS[f \& 3]!}
(line~5). Rule \ruleName{Mask} (\nI{Arithmetic steps}) gives
$0\le f\mathbin{\&}3\le 3$, leaving $\len{\mathit{KINDS}}=4$ to be
established at \code{kindName}'s entry.

\entry{Verification principle}
A module-level representation invariant establishing
$\len{\mathit{KINDS}}=4$ at the entry of every function in the module,
checked to be preserved by every function in the module. Contrast this
with rule \ruleName{Confined} ($\modstate{M}$) in \emph{An Effect
Language for TScert}~\cite{el}: $\modstate{M}$ frames \emph{external}
invariants across mutations of module-private state, whereas P15
requires the dual \emph{internal} invariant stating the length of the
module binding itself at entry.

\entry{Verifier status}
\statusc{Open extension} \lean{SinglePass} verifies the same array
literal when declared locally inside \code{kindName}, but contributes no
length invariant for a module-level \code{const} array at function
entry.

\needspace{18\baselineskip}
\subsection{P16: Cache Location Initialised at First Use (4 sites)}\label{sec:p16}

\begin{lstlisting}[language=TypeScript]
let bmp: Uint16Array | null = null;

/** @requires safeInteger(cp) && 0 <= cp && cp < 65536 */
export function bitsAt(cp: number): number {
  const table = (bmp ??= new Uint16Array(65536));
  return table[cp]!;
}
\end{lstlisting}

\entry{Obligation}
Prove $\mathit{cp}<\len{\mathit{table}}$ at \code{table[cp]!} (line~6).
On the branch where the module \code{let}-cell \code{bmp} already holds
a \code{Uint16Array}, its length must satisfy $\len{\mathit{bmp}}=65536$
so the merge after \code{??=} preserves $\len{\mathit{table}}=65536$.

\entry{Verification principle}
A module-cell representation invariant establishing
$\mathit{bmp}\doteq\code{null}\vee\len{\mathit{bmp}}=65536$ at every
function entry, discharged at initialization and at every write to the
cell. As in P15, rule \ruleName{Confined} ($\modstate{M}$) in \emph{An
Effect Language for TScert}~\cite{el} frames \emph{external} invariants
across mutations of private module \code{let}-cells, whereas P16
requires an \emph{internal} representation invariant constraining the
values stored in the cell itself.

\entry{Verifier status}
\statusc{Open extension} \lean{SinglePass} verifies allocating the
\code{Uint16Array} directly inside \code{bitsAt}, but carries no
cross-call invariant on module-level \code{let}-cells.

\needspace{26\baselineskip}
\subsection{P20: Type Invariant of a Value Read Back from a \texorpdfstring{\code{Map}}{Map} (1 site)}\label{sec:p20}

\begin{lstlisting}[language=TypeScript]
/** @invariant this.index.length === 4352 */
export type ClassTable = { readonly index: Uint16Array };

function build(): ClassTable {
  const index = new Uint16Array(4352);
  return { index };
}

const cache = new Map<string, ClassTable>();

function tableOf(name: string): ClassTable {
  let t = cache.get(name);
  if (t === undefined) {
    t = build();
    cache.set(name, t);
  }
  return t;
}

/** @requires safeInteger(c) && 0 <= c && c < 1114112 */
export function classOf(name: string, c: number): number {
  return tableOf(name).index[c >> 8]!;
}
\end{lstlisting}

\entry{Obligation}
Establish \code{ClassTable}'s type invariant
$\len{t.\mathit{index}}=4352$ at \code{tableOf}'s return (line~17) on
the cache-hit branch where \code{cache.get(name)} returns an existing
entry, enabling \code{index[c >{}> 8]!} in \code{classOf}.

\entry{Verification principle}
A collection-invariant rule paired with a reviewed \code{Map} contract:
reading from a \code{Map} yields a value satisfying a type invariant
whenever every \code{set} into that map stores a value satisfying the
invariant. In \emph{An Effect Language for TScert}~\cite{el}
(Section~6.4), rule \ruleName{Confined} ($\modstate{M}$) solves caching
pattern W5 by framing \emph{external} caller and walker invariants
across private \code{Map}/\code{WeakMap} caches; P20 is the dual
obligation, establishing an \emph{internal} representation invariant on
the cached entries returned by lookup.

\entry{Verifier status}
\statusc{Open extension} \lean{SinglePass} verifies returning a table
directly from its builder (\code{build().index[c >{}> 8]!}), whereas
native \code{Map.prototype.get} carries no type invariant on retrieved
values.

\needspace{30\baselineskip}
% ======================================================================
\section{Arithmetic, Indexing, and Quantified Relations}\label{sec:arith}
% ======================================================================

Seven patterns (34 sites) concern richer term, assertion, or step rules:
non-linear products (P9), shared abstract counts (P10), string indexing
(P11), guarded parallel-array quantification (P12), postfix-counter
prefix fills (P13), term-indexed paths (P17), and running sums of
element lengths (P19). Two of these (P12 and P17) verify in
\lean{SinglePass}.

\needspace{18\baselineskip}
\subsection{P9: Product of Two Variables in an Index (7 sites)}\label{sec:p9}

\begin{lstlisting}[language=TypeScript]
/** @requires safeInteger(width) && 0 < width && width <= 64
 *            && safeInteger(states) && 0 < states && states <= 1024
 *  @requires rows.length === states * width */
export function badProduct(
  rows: Uint16Array, width: number, states: number, state: number,
): number {
  if (!(state >= 0 && state < states)) return 0;
  return rows[state * width]!;
}
\end{lstlisting}

\entry{Obligation}
State $\len{\mathit{rows}}=\mathit{states}\cdot\mathit{width}$ in the
precondition (line~3) and deduce
$0\le\mathit{state}\cdot\mathit{width}<\len{\mathit{rows}}$ at
\code{rows[state * width]!} (line~8).

\entry{Verification principle}
Non-linear arithmetic clauses admitting variable products together with
the entailment
\[
  \len{\mathit{rows}} = \mathit{states}\cdot\mathit{width}
  \;\wedge\; 0\le\mathit{state}<\mathit{states}
  \;\wedge\; 0\le c<\mathit{width}
  \;\vdash\; \mathit{state}\cdot\mathit{width} + c < \len{\mathit{rows}}.
\]

\entry{Verifier status}
\statusc{Open extension} \lean{SinglePass} verifies constant-stride
indexing (such as \code{rows[last * 28 + 27]!} under
\code{rows.length === 6244} and $0\le\mathit{last}<223$), while clause
syntax and arithmetic rules remain linear in variables.

\needspace{22\baselineskip}
\subsection{P10: Equal Counts Computed by Two Functions of One Input (7 sites)}\label{sec:p10}

\begin{lstlisting}[language=TypeScript]
function ends(text: string): number[] {
  const out: number[] = [];
  for (let i = 0; i < text.length; i++) {
    if (text.charCodeAt(i) !== 32) out.push(i);
  }
  return out;
}

function widthsOf(text: string): number[] { /* same loop, pushing 1 */ }

export function paired(text: string): number {
  const e = ends(text);
  const w = widthsOf(text);
  let s = 0;
  for (let k = 0; k < e.length; k++) s += w[k]!;
  return s;
}
\end{lstlisting}

\entry{Obligation}
Prove $k<\len{w}$ at \code{w[k]!} (line~15) from $k<\len{e}$ by
establishing $\len{e}=\len{w}$, even though neither builder's output
length is a linear term in $\len{\mathit{text}}$.

\entry{Verification principle}
A specification term (such as a pure recursive or uninterpreted count
function over strings) naming the count of non-space code units in
\code{text}, so both builders' postconditions equate their result
lengths to that term and the caller derives $\len{e}=\len{w}$.

\entry{Verifier status}
\statusc{Open extension} \lean{SinglePass} verifies reading a single
builder's result below its own length, while the clause language has no
shared abstract count terms.

\needspace{16\baselineskip}
\subsection{P11: String Element Read Below Its Length (5 sites)}\label{sec:p11}

\begin{lstlisting}[language=TypeScript]
export function badElement(text: string): string {
  let out = "";
  for (let i = 0; i < text.length; i++) out = text[i]!;
  return out;
}
\end{lstlisting}

\entry{Obligation}
Type the indexed string read \code{text[i]!} (line~3) as \code{string}
under $0\le i<\len{\mathit{text}}$.

\entry{Verification principle}
Rule \ruleName{Str-Read} of \emph{Tolerant frames}~\cite{tf}: under
$0\le i<\len{s}$ (with $i<2^{32}-1$), \code{y = s[i]} yields a string of
length~$1$.

\entry{Verifier status}
\statusc{Model only / open in SinglePass} Rule \ruleName{Str-Read} is
proved in Lean (\lean{Derivation.stringRead} and
\lean{ArithmeticRules.strRead\_sound}), while \lean{SinglePass} does not
yet emit it (reads of UTF-16 code units via \code{text.charCodeAt(i)}
verify directly).

\needspace{20\baselineskip}
\subsection{P12: Pointwise Relation of Two Arrays (5 sites)}\label{sec:p12}

\begin{lstlisting}[language=TypeScript]
/** @requires safeInteger(i) && 0 <= i && i < fits.length
 *            && fits.length === extras.length
 *  @requires forall k in fits: extras[k] === null
 *            || extras[k].length === fits[k].length */
export function sumWithExtras(
  fits: number[][], extras: (number[] | null)[], i: number,
): number {
  const advances = fits[i]!;
  const e = extras[i] ?? null;
  let s = 0;
  for (let g = 0; g < advances.length; g++) {
    s += advances[g]! + (e === null ? 0 : e[g]!);
  }
  return s;
}
\end{lstlisting}

\entry{Obligation}
Instantiate the guarded parallel-array relation
\[
  \forall k<\len{\mathit{fits}}.\;
    \mathit{extras}[k]\doteq\code{null}
    \;\vee\; \len{\mathit{extras}[k]}=\len{\mathit{fits}[k]}
\]
at index \code{i} so that narrowing \code{e !== null} yields
$\len{e}=\len{\mathit{advances}}$ and proves \code{e[g]!} (line~12) in
bounds. Because \code{extras} has type \code{(number[] | null)[]}, line~9
binds \code{const e = extras[i] ?? null} rather than \code{extras[i]!},
which would falsely assert non-nullness before testing \code{e === null}.

\entry{Verification principle}
Guarded quantified clauses relating two arrays at a shared index,
instantiated at indexed reads of both arrays and released by null
narrowing on the selected element.

\entry{Verifier status}
\status{Verified in SinglePass} Guarded parallel-array quantification
parses, lowers, frames, and instantiates in \lean{SinglePass}
(\lean{IndexShift}, \lean{instantiateEvery\_holds}), verifying
\code{sumWithExtras} directly.

\needspace{22\baselineskip}
\subsection{P13: Typed-Array Prefix Filled with a Postfix Counter (5 sites)}\label{sec:p13}

\begin{lstlisting}[language=TypeScript]
/** @requires safeInteger(n) && 0 <= n && n <= text.length
 *            && ends.length >= n && n <= 2147483647
 *  @ensures safeInteger(result) && 0 <= result && result <= n
 *  @ensures forall i < result of ends:
 *           safeInteger(ends[i]) && 0 < ends[i] && ends[i] <= text.length */
function fill(text: string, ends: Int32Array, n: number): number {
  let count = 0;
  /** @invariant safeInteger(n) && safeInteger(i) && safeInteger(count)
   *             && 1 <= i && i <= n + 1 && count === i - 1
   *             && n <= text.length && n <= ends.length && n <= 2147483647
   *  @invariant forall k < count of ends: safeInteger(ends[k])
   *             && 0 < ends[k] && ends[k] <= text.length */
  for (let i = 1; i <= n; i++) ends[count++] = i;
  return count;
}
\end{lstlisting}

\noindent
The caller allocates \code{ends = new Int32Array(text.length)}, calls
\code{fill(text, ends, text.length)}, and reads
\code{text.codePointAt(ends[0]! - 1)!}.

\entry{Obligation}
Maintain the quantified typed-array prefix invariant
$\forall k<\mathit{count}.\ 0<\mathit{ends}[k]\le\len{\mathit{text}}$
across the postfix write \code{ends[count++] = i} and loop increment
\code{i++} (line~13), and export it through \code{fill}'s postcondition.

\entry{Verification principle}
Two ingredients are needed: extending a typed-array prefix across a
write at the prefix bound followed by a postfix counter increment (rule
\ruleName{View-Every} of \emph{Tolerant frames}~\cite{tf}), and
separating distinct mutable local \code{let}-cells so that updating the
loop counter cell \code{i.value} (which reaches $n+1$, past
$\len{\mathit{text}}$) does not clobber facts supported by the prefix
bound cell \code{count.value}.

\entry{Verifier status}
\statuso{Partially verified} \lean{SinglePass} verifies postfix-counter
prefix fills \code{ends[count++] = v} over freshly allocated typed-array
views (\lean{PromiseCalls.writtenPrefix}). What remains open in
\code{fill} is preserving a prefix bound held in one mutable
\code{let}-cell (\code{count.value}) over a parameter view across an
update to a second mutable \code{let}-cell (\code{i.value}, which shares
the cell key \code{"value"} and exceeds \code{text.length} on loop exit)
without cell-to-cell separation.

\needspace{18\baselineskip}
\subsection{P17: Index into an Element Selected by Another Index (3 sites)}\label{sec:p17}

\begin{lstlisting}[language=TypeScript]
type Cursor = { segment: number; grapheme: number };

/** @requires safeInteger(c.segment) && 0 <= c.segment && c.segment < fits.length
 *            && safeInteger(c.grapheme) && 0 <= c.grapheme
 *            && c.grapheme < fits[c.segment].length */
export function at(fits: number[][], c: Cursor): number {
  const advances = fits[c.segment]!;
  return advances[c.grapheme]!;
}
\end{lstlisting}

\entry{Obligation}
Prove $c.\mathit{grapheme}<\len{\mathit{advances}}$ at
\code{advances[c.grapheme]!} (line~8) from the precondition
\code{c.grapheme < fits[c.segment].length} after reading
\code{const advances = fits[c.segment]!}.

\entry{Verification principle}
Term-indexed paths $\pi[t]$ in clauses and heap facts, paired with a
read link connecting the local \code{advances} to the element selected
at \code{fits[c.segment]} so that the length bound transfers directly to
\code{advances}.

\entry{Verifier status}
\status{Verified in SinglePass} \lean{SinglePass} parses, lowers, frames
(\lean{atIndexFollows\_sound}), and instantiates
(\lean{instantiateAtIndex\_holds}) indexed-element facts via
\lean{IndexFact.atIndex}, verifying \code{at} directly.

\needspace{20\baselineskip}
\subsection{P19: Running Sum of Element Lengths (2 sites)}\label{sec:p19}

\begin{lstlisting}[language=TypeScript]
/** @requires segments.length <= 1000
 *            && forall i in segments: segments[i].length <= 1000 */
export function sources(segments: string[]): number {
  let total = 0;
  /** @invariant safeInteger(i) && 0 <= i && safeInteger(total)
   *             && 0 <= total && total <= 1000 * i */
  for (let i = 0; i < segments.length; i++) total += segments[i]!.length;
  const offsets = new Int32Array(total + 1);
  let s = 0;
  for (let i = 0, at = 0; i < segments.length; at += segments[i++]!.length) {
    s += offsets[at]!;
  }
  return s;
}
\end{lstlisting}

\entry{Obligation}
Prove $0\le\mathit{at}<\len{\mathit{offsets}}$ at \code{offsets[at]!}
(line~11), where \code{at} accumulates non-constant element lengths
$\len{\mathit{segments}[i]}$ across the second loop.

\entry{Verification principle}
An aggregate summation term over a slice of element lengths so that the
loop invariant
\[
  \mathit{at} + \sum_{i\le k<\len{\mathit{segments}}} \len{\mathit{segments}[k]}
  \;=\; \mathit{total}
\]
can be stated and maintained on each iteration.

\entry{Verifier status}
\statusc{Open extension} \lean{SinglePass} verifies indexing
\code{offsets} directly at \code{total}, while neither \lean{SinglePass}
nor the companion papers provide aggregate element-length sum terms.

\needspace{26\baselineskip}
% ======================================================================
\section{External Specifications}\label{sec:external}
% ======================================================================

Two patterns (20 direct sites, plus 33 dependent sites) depend on
properties of code or data outside the analysed TypeScript source:
offline-generated lookup tables (P5) and string-library postconditions
(P14).

\needspace{22\baselineskip}
\subsection{P5: Bounds Held by the Contents of Generated Tables (16 sites)}\label{sec:p5}

\begin{lstlisting}[language=TypeScript]
/** @ensures result.length === 4 */
function decode(): Uint8Array {
  const packed = "\u0003\u0001\u0002\u0000";
  const out = new Uint8Array(4);
  for (let i = 0; i < 4; i++) out[i] = packed.charCodeAt(i) & 0xff;
  return out;
}

/** @requires safeInteger(state) && 0 <= state && state < 4 */
export function next(state: number): number {
  const t = decode();
  return t[t[state]!]!;
}
\end{lstlisting}

\entry{Obligation}
Prove $0\le t[\mathit{state}]<\len{t}$ at the outer read
\code{t[t[state]!]!} (line~12). Because \code{decode}'s postcondition
states only $\len{\mathit{result}}=4$ and its loop copies bytes from the
packed string literal without checking their values, no fact bounds the
table's contents.

\entry{Verification principle}
Establish $\forall i<\len{t}.\ 0\le t[i]<\len{t}$ at \code{decode}'s
return, either by statically evaluating the constant string literal
during verification, by a runtime validation loop inside \code{decode},
or as an explicit external assumption on the offline table generator.

\entry{Verifier status}
\statusc{Open extension / trust boundary} \lean{SinglePass} verifies the
inner read \code{t[state]!} against \code{decode}'s stated length, while
constant-table evaluation is not part of the verifier.

\needspace{18\baselineskip}
\subsection{P14: Length Contracts of String and Array Natives (4 sites, 33 dependent)}\label{sec:p14}

\begin{lstlisting}[language=TypeScript]
export function bad(text: string): number {
  const normalized = text.replace(/ +/g, " ");
  const breaks = new Uint8Array(text.length + 1);
  let s = 0;
  /** @invariant safeInteger(i) && 0 <= i && normalized.length <= text.length
   *             && breaks.length === text.length + 1 */
  for (let i = 0; i < normalized.length; i++) s += breaks[i]!;
  return s;
}
\end{lstlisting}

\entry{Obligation}
Establish $\len{\mathit{normalized}}\le\len{\mathit{text}}$ at loop
entry (line~5) from \code{text.replace(/ +/g, " ")} so \code{breaks[i]!}
(line~7) is in bounds.

\entry{Verification principle}
Reviewed host length postconditions for the string and array natives
used at these sites---for example, $\len{r}\le\len{s}$ for
\code{r = s.replace(e, t)} whenever every match of regex \code{e} has
length at least $\len{t}$.

\entry{Verifier status}
\statusc{Open extension / host contract} \lean{SinglePass} verifies
methods with exact reviewed length contracts such as
\code{text.slice(0, k)} (length $k$ for $0\le k\le\len{\mathit{text}}$),
while regex-dependent \code{replace} length contracts are not specified.

\clearpage
% ======================================================================
\section{Summary and Trust Boundary}\label{sec:summary}
% ======================================================================

\subsection{Master Summary of Catalogued Patterns}

Table~\ref{tab:summary} summarizes all 22 patterns across the baseline
215-site blocked corpus: the four flagship patterns of
Section~\ref{sec:verified} followed by the 18 entries of
Sections~\ref{sec:framing}--\ref{sec:external} in decreasing order of
Pretext site count, together with P22.

\begin{table}[!htbp]
\centering\footnotesize
\renewcommand{\arraystretch}{1.04}
\begin{tabularx}{\textwidth}{@{}l r >{\raggedright\arraybackslash}p{0.29\textwidth} >{\raggedright\arraybackslash}X@{}}
\toprule
Pattern & Sites & Proof obligation & Theory and verifier status\\
\midrule
P1 (\S\ref{sec:verified}) & 50 & Lower bound across visitor & \status{Verified in SinglePass} (\ruleName{(Grow)}~\cite{tf}); 48 walker sites also require P22.\\
P2 (\S\ref{sec:verified}) & 31 & Growth by known count & \status{Verified in SinglePass} (two-state \code{@ensures} + regional spine frame).\\
P8 (\S\ref{sec:verified}) & 10 & Write at array length & \status{Verified in SinglePass} (\ruleName{Write-End}~\cite{tf} + \lean{denseAppendAtoms\_sound}).\\
P21 (\S\ref{sec:verified}) & 1\rlap{\textsuperscript{\ddag}} & Shift-left / bitwise-OR index & \status{Verified in SinglePass} (\ruleName{Shr}, \ruleName{Shl}, \ruleName{Mask}, \ruleName{Or}~\cite{tf}); sits behind P5.\\
\midrule
P3 (\S\ref{sec:p3}) & 19 & Quantified fact across field write & \status{Verified in SinglePass} (\lean{Region.pathCells} $\cap$ fields $=\emptyset$, \lean{Survival.everyCells}).\\
P4 (\S\ref{sec:p4}) & 18 & Fresh path separation; \code{i + 1 < n} & \statuso{Partial}: \code{paths} verified (\lean{PathRegionsApart}, \lean{restateSpines}); \code{straight} open.\\
P5 (\S\ref{sec:p5}) & 16 & Value bounds of generated table & \statusc{Open}: literal evaluation, runtime check, or generator assumption.\\
P6 (\S\ref{sec:p6}) & 12 & Quantified fact across field call & \statuso{Model + prototype}: parameter fields verified in \lean{SinglePass}; dynamic fields via \ruleName{Fields-Disj}~\cite{el}.\\
P7 (\S\ref{sec:p7}) & 10 & Captured facts in local closures & \statuso{Model + prototype} for visitors (\ruleName{Callback-Inv}~\cite{el}); capture \code{@requires} open.\\
P9 (\S\ref{sec:p9}) & 7 & Product of two variables & \statusc{Open} (constant strides verified in \lean{SinglePass}).\\
P10 (\S\ref{sec:p10}) & 7 & Shared count across two builders & \statusc{Open} (needs pure/uninterpreted count terms).\\
P11 (\S\ref{sec:p11}) & 5 & String element read below length & \statusc{Model only}: \ruleName{Str-Read}~\cite{tf} proved in Lean; not emitted by \lean{SinglePass}.\\
P12 (\S\ref{sec:p12}) & 5 & Guarded two-array quantification & \status{Verified in SinglePass} (\lean{IndexShift}, \lean{instantiateEvery\_holds}; with \code{extras[i] ?? null}).\\
P13 (\S\ref{sec:p13}) & 5 & Postfix-counter prefix fill & \statuso{Partial}: fresh-view fill verified (\lean{writtenPrefix}); two-cell \code{fill} open.\\
P14 (\S\ref{sec:p14}) & 4 & Length contract of \code{replace} & \statusc{Open}: reviewed host postcondition (33 dependent sites).\\
P15 (\S\ref{sec:p15}) & 4 & Module constant array length & \statusc{Open}: internal module invariant (dual to \ruleName{Confined}~\cite{el}).\\
P16 (\S\ref{sec:p16}) & 4 & Module lazy cache cell invariant & \statusc{Open}: internal cell invariant (dual to \ruleName{Confined}~\cite{el}).\\
P17 (\S\ref{sec:p17}) & 3 & Term-indexed path in clauses & \status{Verified in SinglePass} (\lean{IndexFact.atIndex}, \lean{instantiateAtIndex\_holds}).\\
P18 (\S\ref{sec:p18}) & 2 & Two-field \code{@stable} guard & \statusc{Model only}: \ruleName{(Rigid)}~\cite{tf}; \lean{SinglePass} keeps direct sibling writes only.\\
P19 (\S\ref{sec:p19}) & 2 & Running sum of element lengths & \statusc{Open} (needs aggregate length-sum terms).\\
P20 (\S\ref{sec:p20}) & 1 & Type invariant of \code{Map} entry & \statusc{Open}: internal map-entry invariant (dual to \ruleName{Confined} / W5~\cite{el}).\\
P22 (\S\ref{sec:p22}) & 1\rlap{\textsuperscript{\dag}} & Local cursor \& handle in callback & \statuso{Model + prototype}: 9 rules in~\cite{el} (\lean{model\_rules}) \statusc{conj.}; 105/105 Pretext units pass with 0 \code{paperAssume}s.\\
\bottomrule
\end{tabularx}
\caption{Master summary of Pretext's 22 bound patterns (baseline audit corpus \code{9268edb1}: 118 sites verified in \lean{SinglePass} + 97 partial/model/open = 215 blocked sites). \textsuperscript{\ddag}P21's 1 site sits behind P5 and is counted in P5's 16 sites. \textsuperscript{\dag}P22 is 1 call-site pattern acting as the second blocker on 48 of P1's walker sites.}\label{tab:summary}
\end{table}

\paragraph{Sites in no pattern.}
Three bound rewrite sites of the baseline corpus (\code{9268edb1}) fall
outside P1--P21:
(i)~indexing a parallel normalized start-offset list
(\code{scanStarts[k]} in \code{analyzeText}),
(ii)~updating a sentinel-bounded break array at a shifted offset
(\code{sourceBreaks[after]} in \code{analyzeText}), and
(iii)~indexing a pair-strided bracket stack at $2k$ and $2k+1$
(\code{openBrackets[opened + 1]} in \code{readScriptRuns}).
The 422-site corpus also excludes 34 rewrite markers whose operand is
not an index read.

\FloatBarrier
\subsection{Trust Boundary: Facts About Unanalysed Code}\label{sec:trust}

Four obligations across the catalogue depend on code or data outside the
analysed TypeScript files and therefore constitute explicit trust
assumptions rather than internal proofs:
\begin{enumerate}[leftmargin=1.5em,itemsep=2pt]
\item \textbf{Library length contracts (P14):} Postconditions on
  built-in methods such as \code{String.prototype.replace} are reviewed
  specifications of host code.
\item \textbf{Generated tables (P5):} Value bounds on packed lookup
  tables trust the offline generator unless established by constant
  evaluation or a runtime check at decode time.
\item \textbf{Unanalysed callbacks (P1, P22):} An effect bound stated on
  a callback parameter of an exported API function (whether closed as in
  P1 or polymorphic as in~\cite{el}) is discharged for analysed callers
  and assumed of external callers.
\item \textbf{Exported preconditions and handle invariants (P22):}
  Handles passed into exported functions from unanalysed clients are
  assumed to satisfy their declared \code{@invariant}, as with any
  precondition on an exported entry point (whereas caller-supplied
  streaming cursors may be validated by explicit boundary checks).
\end{enumerate}

\begin{thebibliography}{4}\small\setlength{\itemsep}{1pt}\setlength{\parskip}{0pt}
\bibitem{tf} TScert Working Note. \emph{Tolerant frames: a modal account of fact preservation without ownership}. 27 November 2026.
\bibitem{el} TScert Working Paper. \emph{An Effect Language for TScert: Relational Effects, Bounded Effect Polymorphism, and Tolerant Framing}. 27 November 2026.
\bibitem{pretext} Pretext. TypeScript text-layout library (\texttt{chenglou/pretext}), baseline bound-rewrite audit corpus at commit \code{9268edb1} and walker/builder verification against upstream commit \code{0dafecb3}.
\end{thebibliography}

\end{document}
